\subsection{互奇异测度}

对 T 上的两个实测度 \(\rho,\sigma\) ，回忆一下，称 \(\rho\) 和 \(\sigma\) 是互奇异的，如果 \(\inf(|\rho|,|\sigma|)=0\) 在 \(\mathcal{M}(\mathrm{T})\) 中成立（Ch. II, §1, No.~1）。已知与一个给定测度互奇异的诸实测度构成一个带（Ch. II, §1, No.~5, 定理 1）。这个定义可以立即推广到复测度的情形。

定义 4. --- 称 T 上的一个复测度 \(\theta\) 集中于 T 的一个子集 M 上，或称 M 承载 \(\theta\) ，如果 \(\mathbb{C}M\) 对 \(\theta\) 是局部可忽略的。

集 M 承载 \(\theta\) 当且仅当它承载 \(|\theta|\) 。说 M 承载 \(\theta\) ，等价于说 M 是 \(\theta\) -可测的且 \(\theta = \varphi_{M} \cdot \theta\) 。若 \(\theta\) 集中于 M，则每个以 \(|\theta|\) 为基的测度都集中于 M。

命题 13. --- 为使 T 上的两个复测度 \(\rho\) 和 \(\sigma\) 互奇异，必须且只须在 T 中存在两个不相交的集 R 和 S ，使得 \(\rho\) 集中于 R 且 \(\sigma\) 集中于 S；可以取 R 和 S 为普遍可测的。

置 \(\mu = |\rho|\) ，\(\nu = |\sigma|\) ，\(\lambda = \mu + \nu\) ；由于 \(\mu\) 和 \(\nu\) 被 \(\lambda\) 优超，存在两个局部 \(\lambda\) -可积函数 u 和 v （依 §3, No.~4, 命题 7 可设它们是普遍可测的），使得 \(\mu = u \cdot \lambda\) ，\(\nu = v \cdot \lambda\) 。于是

\[
\inf (| \rho |, | \sigma |) = \inf (\mu , \nu) = \inf (u, v) \cdot \lambda
\]

（No.~2, 命题 2 的推论）。设 A （相应地，B ）是由 \(t \in \mathbf{T}\) （使 \(u(t) > 0\) 且 \(v(t) = 0\) （相应地，\(u(t) = 0\) 且 \(v(t) > 0\) ）者）组成的集。若 \(\rho\) 和 \(\sigma\) 互奇异，则 \(\inf(u, v) = 0\) 局部 \(\lambda\) -几乎处处成立（No.~3, 命题 3 的推论 2），于是不相交的普遍可测集 A 和 B 分别承载 \(\mu\) 和 \(\nu\) 。反之，设 \(\mu\) 和 \(\nu\) 分别由不相交的集 R 和 S 承载；\(\varphi_{\mathrm{R}}\) 对测度 \(\mu = u \cdot \lambda\) 可测，且 \(\mu = \varphi_{\mathrm{R}} \cdot \mu\) 。由 No.~4 的命题 8，函数 \(u' = u\varphi_{\mathrm{R}}\) 是 \(\lambda\) -可测的，且 \(\mu = u' \cdot \lambda\) 。类似地，若 \(v' = v\varphi_{\mathrm{S}}\) ，则 \(\nu = v' \cdot \lambda\) ；注意到 \(\inf(u', v') = 0\) （No.~2, 命题 2 的推论），即得结论。

推论 1. --- 对 T 上的每个实测度 \(\nu\) ，存在两个不相交的集 M,N ，分别承载 \(\nu^{+}\) 和 \(\nu^{-}\) 。

\pandocbounded{\includegraphics[keepaspectratio]{images/part002_61ad429f.jpg}}

必须注意不要把测度 \(\nu\) 的支集的概念与 \(\nu\) 所集中的集的概念相混淆。\(\nu\) 的支集 S 是承载 \(\nu\) 的最小闭集（Ch. III, §2, No.~2, 命题 2 及 Ch. IV, §2, No.~2, 命题 5）。然而，可能存在 S 的与 S 不同却承载 \(\nu\) 的子集。更一般地，可以有 \(\inf(\mu,\nu)=0\) ，对两个具有相同支集的正测度 \(\mu\) 和 \(\nu\) 而言（习题 5）。

还要注意，承载 \(\nu\) 的诸集之交是由 \(t \in T\) （使 \(|\nu|(\{t\}) > 0\) 者）组成的点的集，它可以是空集（例如 Lebesgue 测度的情形）；因此，一般不存在承载 \(\nu\) 的最小集。

推论 2. --- 设 \(\rho\) 和 \(\sigma\) 是两个互奇异的复测度，\(\rho'\) 和 \(\sigma'\) 是两个复测度，分别具有关于 \(\rho\) 和 \(\sigma\) 的密度；则 \(\rho'\) 和 \(\sigma'\) 互奇异。

推论 3. --- 设 \(\rho\) 和 \(\sigma\) 是两个互奇异的复测度；则 \(|\rho + \sigma| = |\rho| + |\sigma|\) 。

用 v （相应地，w ）表示一个绝对值为 1 的普遍可测函数，使得 \(\rho = v \cdot |\rho|\) （相应地，\(\sigma = w \cdot |\sigma|\) ）（定理 2 的推论 3），并用 A 表示一个承载 \(\rho\) 的普遍可测集，使得 B = C A 承载 \(\sigma\) （命题 13）；则也有 \(\rho + \sigma = (v\varphi_{\mathrm{A}} + w\varphi_{\mathrm{B}}) \cdot (|\rho| + |\sigma|)\) 。由于函数 \(v\varphi_{A} + w\varphi_{B}\) 的绝对值等于 1 ，推论由命题 2 得出。

定理 3（Lebesgue）. --- T 上的每个复测度 \(\theta\) 都可以以一种且仅一种方式写成 \(\theta = g \cdot \mu + \theta'\) 的形式，其中 \(g\) 是局部 \(\mu\) -可积的，\(\theta'\) 是一个与 \(\mu\) 互奇异的测度。此时 \(|\theta| = |g| \cdot \mu + |\theta'|\) 。

当 \(\theta\) 是正的时，这由 F. Riesz 定理（Ch. II, §1, No.~5, 定理 1）（应用于 T 上的实测度所组成的完全格序空间 \(\mathcal{M}(T)\) 以及 \(\mu\) 在此空间中生成的带，并注意到 No.~5 的定理 2 的推论 2）立即得出；此外，此时 \(\theta'\) 和 \(g \cdot \mu\) 都是正的，这蕴涵 \(g\) 局部 \(\mu\) -几乎处处为正（命题 3 的推论 3）。为处理 \(\theta\) 不是正的情形，置 \(\nu = |\theta|\) ，\(\nu = f \cdot \mu + \nu'\) （其中 \(f\) 是正的，且 \(\nu'\) 与 \(\mu\) 互奇异），并置 \(\theta = v \cdot \nu\) ，其中 \(v\) 是一个绝对值为 1 的普遍可测函数（定理 2 的推论 3）。于是我们有（命题 8）\(\theta = g \cdot \mu + \theta'\) ，其中 \(g = vf\) （从而 \(|g| = f\) ）且 \(\theta' = v \cdot \nu'\) （从而依命题 2 有 \(|\theta'| = \nu'\) ）；测度 \(\theta'\) 与 \(\mu\) 依命题 13 的推论 2 互奇异。只须再建立分解的唯一性。于是，设 \(\theta = g \cdot \mu + \theta' = g_1 \cdot \mu + \theta'_1\) ，其中 \(\theta'\) 和 \(\theta'_1\) 都与 \(\mu\) 互奇异；\(|\theta' - \theta'_1|\) 被 \(|\theta'| + |\theta'_1|\) 优超，因此 \(\theta' - \theta'_1\) 与 \(\mu\) 互奇异，从而也与 \((g_1 - g) \cdot \mu\) 互奇异。于是关系 \(\theta' - \theta'_1 = (g_1 - g) \cdot \mu\) 蕴涵两个成员都为零，这就证明了唯一性。

回忆（No.~5, 定理 2 及概说），空间 \(\mathrm{L}_{\mathrm{loc}}^{1}(\mathrm{T},\mu;\mathbf{C})\) 可以（借助映射 \(g\mapsto g\cdot\mu\) ）等同于 \(\mathcal{M}_{\mathbf{C}}(\mathrm{T})\) 的一个子空间。在这一约定下，定理 3 取下列形式：

推论. --- 存在一个投影子 \(p\) （由空间 \(\mathcal{M}_{\mathbf{C}}(\mathrm{T})\) 到空间 \(\mathrm{L}_{\mathrm{loc}}^{1}(\mathrm{T},\mu ;\mathbf{C})\) 上的），其核 \(\overline{p}^{1}(0)\) 是与 \(\mu\) 互奇异的复测度所组成的集，使得

\[
| \theta | = | p (\theta) | + | \theta - p (\theta) |, \quad p (| \theta |) = | p (\theta) |
\]

对每个复测度 \(\theta\) 都成立。

若把 p 限制在有界测度所成之集上，就得到一个投影子 \(p^{1}\) （由空间 \(\mathcal{M}_{\mathbf{C}}^{1}(\mathrm{T})\) 到空间 \(\mathrm{L}_{\mathbf{C}}^{1}(\mathrm{T},\mu)\) 上的）；关系 \(\|\theta\|=|\theta|(1)\) 蕴涵 \(\|\theta\|=\|p^{1}(\theta)\|+\|\theta-p^{1}(\theta)\|\) 对每个有界复测度 \(\theta\) 成立。

